\documentclass[11pt]{amsart}
\usepackage{amsmath,amssymb}
\setlength{\emergencystretch}{3em}
\tolerance=600
\usepackage{tikz-cd}
\usepackage{tikz}
\usepackage[hidelinks]{hyperref}

\newtheorem{thm}{Theorem}[section]
\newtheorem{lem}[thm]{Lemma}
\newtheorem{cor}[thm]{Corollary}
\newtheorem{prop}[thm]{Proposition}
\theoremstyle{definition}
\newtheorem{defi}[thm]{Definition}
\newtheorem{rem}[thm]{Remark}
\newtheorem{ex}[thm]{Example}
\newtheorem{que}[thm]{Question}

\newcommand{\Coh}{\mathbf{Coh}}
\newcommand{\Heyt}{\mathbf{Heyt}}
\newcommand{\Mod}{\mathrm{Mod}}
\newcommand{\Set}{\mathbf{Set}}
\newcommand{\Sub}{\mathrm{Sub}}
\newcommand{\FF}{\mathsf F}
\newcommand{\EE}{\mathsf E}
\newcommand{\PP}{\mathsf P}
\newcommand{\BB}{\mathsf B}
\newcommand{\HH}{\mathsf H}
\newcommand{\KK}{\mathsf K}
\newcommand{\cC}{\mathcal C}
\newcommand{\cD}{\mathcal D}
\newcommand{\cK}{\mathcal K}

\begin{document}

\title[The free Heyting category on a pretopos]{The free Heyting category on a pretopos is a pretopos only in the Boolean case: the intuitionistic translation and the pretopos completion}
\author{Ryota Shibusawa}
\thanks{Corresponding author. E-mail: \texttt{r-shibusawa@daiichi-koudai.ac.jp}}
\address{Daiichi Institute of Technology}
\email{r-shibusawa@daiichi-koudai.ac.jp}
\keywords{Heyting category, pretopos completion, coherent category, intuitionistic first-order logic, Kripke model, imaginaries, Morleyization, Boolean category}
\subjclass[2020]{03G30 (primary), 18F10, 03F55, 03C90 (secondary)}

\begin{abstract}
Reading a coherent theory in intuitionistic first-order logic is the
free Heyting category $\HH\cC$ on its coherent syntactic category
$\cC$; adding the imaginaries is the pretopos completion $\PP$. We
prove that the two operations commute only trivially: the canonical
functor $\HH\PP\cC\to\PP\HH\cC$ is always full and faithful, and it is
an equivalence if and only if $\cC$ is Boolean, in which case $\HH$
does nothing. The objects we show to be
missing from $\HH\PP\cC$ are two copies of an object glued along the intuitionistic negation of a
non-complemented subobject; they are excluded by an invariant of the
objects of $\HH\PP\cC$ in Kripke models over a fork, whose existence
for every non-Boolean $\cC$ is a form of the homomorphism preservation
theorem. $\HH$ commutes with the positive completion, and with the exact
completion for propositional theories but not in general. For subobjects there is a preservation theorem: an
intuitionistic formula is coherent if and only if its forcing is
reflected along the transitions which are isomorphisms of models; so
the equivalence relations on coherent imaginaries with a quotient in
$\HH\PP\cC$ are exactly the coherent ones. For a propositional theory the image of the comparison is
characterised by a condition on zigzags in the finite Kripke models of
its finite subtheories, and for a finite one in its own models. The Morleyization being
$\HH$ followed by excluded middle, which commutes with $\PP$, its
Galois criterion says when excluded middle repairs the comparison.
\end{abstract}

\maketitle

\section{Introduction}

A coherent theory $T$ has a coherent syntactic category $\cC_T$, and
three constructions on it are the subject of this paper and of its
predecessor \cite{Morley}. The pretopos completion $\PP\cC$ adds finite
coproducts of sorts and quotients by definable equivalence relations,
the coherent imaginaries \cite[A3.3, D1.4]{Elephant}. The Morleyization
$\BB\cC$ is the syntactic category of $T$ read in classical first-order
logic. Between them lies the Heyting completion $\HH\cC$, the syntactic
category of $T$ read in intuitionistic first-order logic, and
$\BB\simeq\KK\HH$ with $\KK$ the addition of excluded middle
\cite[\S9]{Morley}.

In \cite{Morley} the order of translation and completion was studied
for $\BB$: there is a full and faithful comparison
$\BB\PP\cC\to\PP\BB\cC$, which is an equivalence exactly when a Galois
criterion holds \cite[Thm.~6.1]{Morley}; it fails for the theory of
complete multipartite graphs and holds, for example, for the theory of
an equivalence relation. It was also shown that $\KK$ commutes with
$\PP$ on all Heyting categories \cite[Thm.~9.2]{Morley}, so that the
comparison for $\BB$ is the image under $\KK$ of the comparison
\[
c_\cC\colon\HH\PP\cC\longrightarrow\PP\HH\cC
\]
for $\HH$. The question left open was when $c_\cC$ itself is an
equivalence. The answer is the following.

\begin{thm}\label{thm:main}
For a small coherent category $\cC$ the functor $c_\cC$ is full and
faithful, and the following are equivalent.
\begin{enumerate}
\item[(i)] $c_\cC\colon\HH\PP\cC\to\PP\HH\cC$ is an equivalence.
\item[(ii)] $\HH\PP\cC$ is a pretopos.
\item[(iii)] $\cC$ is Boolean.
\end{enumerate}
If they hold, the units $\cC\to\HH\cC$ and $\PP\cC\to\HH\PP\cC$ are
equivalences.
\end{thm}

For a pretopos $\cC$ this says that the free Heyting category on $\cC$
is a pretopos if and only if $\cC$ is Boolean
(Corollary~\ref{cor:pretopos}).

\subsection*{What is new and what is known}
New, to the author's knowledge: Theorem~\ref{thm:main}; the invariant
that proves it (Lemma~\ref{lem:rigid}); the commutation of $\HH$ with
the positive completion (Proposition~\ref{prop:F}) and the resulting
location of the failure (Corollary~\ref{cor:E},
Proposition~\ref{prop:E}); the preservation theorem for rigid formulas
and the description of the effective equivalence relations
(Theorems~\ref{thm:rigid} and~\ref{thm:effective}); and the reading of the
Galois criterion of \cite{Morley} as the statement that $c_\cC$ is
inverted by $\KK$ (Corollary~\ref{cor:K}). Known and used: syntactic
categories of first-order theories are Heyting categories and classify
their models in Heyting categories \cite[D1.4]{Elephant}; the pretopos
completion \cite[A3.3]{Elephant}; completeness of coherent logic for
set-valued models \cite[D1.5]{Elephant}, \cite{MR77}; the homomorphism
preservation theorem of \L o\'s, Tarski and Lyndon, of which
Lemma~\ref{lem:cross} is the form needed here and is proved in full;
the Heyting structure of presheaf categories, that is, Kripke
semantics, and Joyal's completeness theorem for it \cite{MR95}; and from \cite{Morley} the facts that $\PP\cD$ is Heyting
when $\cD$ is, with the expected universal property, and that $\KK$
commutes with $\PP$.

\subsection*{Conventions}
$\Coh$ is the $(2,1)$-category of small coherent categories, coherent
functors and natural isomorphisms, and $\Heyt$ that of small Heyting
categories (coherent categories in which every pullback functor
$f^*\colon\Sub(B)\to\Sub(A)$ has a right adjoint $\forall_f$), Heyting
functors (coherent functors preserving the $\forall_f$) and natural
isomorphisms \cite[A1.4]{Elephant}. A Heyting functor preserves the
implication of subobjects and the pseudo-complement $\neg U=(U
\Rightarrow0)$, since $U\Rightarrow V=\forall_i(i^*V)$ for the
inclusion $i\colon U\to A$. A coherent category is \emph{Boolean} if
every subobject is complemented. A model of $\cC$ is a coherent functor
$\cC\to\Set$, and $\Mod(\cC)$ is the category of models with all
natural transformations, the homomorphisms. The universal properties
below are stated for small categories but hold, with the same proofs,
for targets such as $\Set^{\cK}$, $\cK$ a small category; alternatively
one replaces $\Set^{\cK}$ by a small full subcategory closed under the
operations of a Heyting pretopos and containing the objects in
question.

\section{The two completions and the comparison}

\subsection*{The Heyting completion}
For $\cC$ in $\Coh$ let $\HH\cC$ be the syntactic category of the
internal coherent theory of $\cC$ (a sort for every object, a function
symbol for every morphism, the coherent sequents valid in $\cC$) read in
intuitionistic first-order logic with possibly empty sorts (the initial
object is a sort), and $\eta\colon\cC\to\HH\cC$ the
coherent functor sending an object to its sort. $\HH\cC$ is a Heyting
category, and for every Heyting category $\cD$ restriction along $\eta$
is an equivalence of groupoids
\[
\Heyt(\HH\cC,\cD)\;\simeq\;\Coh(\cC,\cD):
\]
a coherent functor $\cC\to\cD$ is a model of the internal theory in
$\cD$, it interprets every first-order formula, intuitionistic logic is
sound in Heyting categories, and an isomorphism of models preserves the
interpretations of all formulas \cite[D1.4]{Elephant},
\cite[\S9]{Morley}. Every object of $\HH\cC$ is a subobject of some
$\eta A$: it is a formula in a context $x_1{:}A_1,\dots,x_n{:}A_n$, a
subobject of $\eta(A_1\times\dots\times A_n)$.

The forgetful $2$-functor $\Heyt\to\Coh$ is not full, a coherent
functor between Heyting categories need not be Heyting. So $\HH$ is a
free construction and not a reflector; the join lemma of
\cite{TheoryComp,Morley}, which concerns reflectors, does not apply to
it, and we compare the composites $\HH\PP$ and $\PP\HH$ directly.

\begin{lem}\label{lem:boolean}
If $\cC$ is Boolean, then $\eta\colon\cC\to\HH\cC$ is an equivalence.
\end{lem}
\begin{proof}
A Boolean coherent category is Heyting, with $\forall_fU=\neg\exists_f
\neg U$. Let $G\colon\cC\to\cD$ be coherent with $\cD$ Heyting. Then
$G$ preserves complements and images, so $G(\forall_fU)$ is a
complement $W$ of $\exists_{Gf}V'$, where $V'=G(\neg U)$ is a
complement of $V=GU$. In a Heyting category this $W$ is $\forall_{Gf}
V$: for a subobject $W'$, $(Gf)^*W'\le V$ iff $(Gf)^*W'\wedge V'=0$
iff $W'\wedge\exists_{Gf}V'=0$ (Frobenius reciprocity) iff $W'\le W$.
So every coherent functor from $\cC$ to a Heyting category is Heyting,
and the identity of $\cC$ has the universal property of $\eta$.
\end{proof}

\subsection*{The pretopos completion of a Heyting category}
We write $\iota\colon\cC\to\PP\cC$ for the unit of the pretopos
completion; it is full and faithful, and restriction along it is an
equivalence between the categories of coherent functors $\PP\cC\to
\mathcal E$ and $\cC\to\mathcal E$ with all natural transformations
(not only the groupoids $\Coh(\PP\cC,\mathcal E)$ and $\Coh(\cC,
\mathcal E)$), when $\mathcal E$ is a
pretopos \cite[A3.3]{Elephant}. In particular $\Mod(\PP\cC)\simeq
\Mod(\cC)$. Every object of $\PP\cC$ is the quotient of a finite
coproduct of objects $\iota A$ by an equivalence relation.

\begin{lem}[{\cite[Lemma~9.1]{Morley}}]\label{lem:PHeyt}
If $\cD$ is a Heyting category, then $\PP\cD$ is a Heyting pretopos and
$\iota\colon\cD\to\PP\cD$ is a Heyting functor; and for a Heyting
pretopos $\mathcal E$ a coherent functor $\PP\cD\to\mathcal E$ is
Heyting if and only if its restriction to $\cD$ is.
\end{lem}

Consequently $\PP\HH\cC$ is the free Heyting pretopos on $\cC$:
$\Heyt(\PP\HH\cC,\mathcal E)\simeq\Coh(\cC,\mathcal E)$ for Heyting
pretoposes $\mathcal E$.

\subsection*{The comparison}
The coherent functor $\PP\eta\colon\PP\cC\to\PP\HH\cC$ has a Heyting
target, so it extends, uniquely up to a unique isomorphism, to a
Heyting functor
\[
\begin{tikzcd}[column sep=large]
\PP\cC\arrow[r,"\eta_{\PP\cC}"]\arrow[dr,"\PP\eta"']&\HH\PP\cC
\arrow[d,"c_\cC"]\\
&\PP\HH\cC.
\end{tikzcd}
\]

\begin{lem}\label{lem:ff}
$c_\cC$ is full and faithful. It is an equivalence if and only if
$\HH\PP\cC$ is a pretopos.
\end{lem}
\begin{proof}
For a Heyting pretopos $\mathcal E$, by Lemma~\ref{lem:PHeyt} and the
universal properties,
\begin{align*}
\Heyt(\PP\HH\PP\cC,\mathcal E)&\simeq\Heyt(\HH\PP\cC,\mathcal E)\simeq
\Coh(\PP\cC,\mathcal E)\\&\simeq\Coh(\cC,\mathcal E)\simeq
\Heyt(\PP\HH\cC,\mathcal E),
\end{align*}
by restriction to $\cC$. So there is a Heyting equivalence $e\colon\PP
\HH\PP\cC\simeq\PP\HH\cC$ compatible with the functors from $\cC$, and
$e\circ\iota_{\HH\PP\cC}$ is a Heyting functor $\HH\PP\cC\to\PP\HH\cC$
whose restriction to $\PP\cC$ is a coherent functor into a pretopos
restricting on $\cC$ as $\PP\eta$ does; hence $e\circ\iota_{\HH\PP\cC}
\cong c_\cC$. The unit $\iota_{\HH\PP\cC}$ is full and faithful, and it
is an equivalence if and only if $\HH\PP\cC$ is a pretopos.
\end{proof}

\section{Kripke models and an invariant}

Let $\cK$ be a small category. The category $\Set^{\cK}$ is a Heyting
pretopos, and limits, images, unions, coproducts and quotients in it
are computed at each object of $\cK$; so a coherent functor $\cC\to
\Set^{\cK}$ is the same as a functor $M\colon\cK\to\Mod(\cC)$, a
\emph{Kripke model} of $\cC$ over $\cK$; the transitions are arbitrary
homomorphisms and need not be injective. We write $M_k$ for the model
at $k$. The pseudo-complement in $\Set^{\cK}$ is the Kripke negation:
for a subfunctor $S\subseteq A$,
\begin{equation}\label{eq:neg}
(\neg S)(k)=\{x\in A(k)\mid A(w)x\notin S(l)\text{ for every arrow }
w\colon k\to l\},
\end{equation}
since the right-hand side is a subfunctor disjoint from $S$, and every
subfunctor $W$ disjoint from $S$ lies in it ($x\in W(k)$ gives
$A(w)x\in W(l)$).

A Kripke model $M$ extends to
\begin{align*}
M^{\HH}&\colon\HH\cC\to\Set^{\cK},& M^{\PP}&\colon\PP\cC\to\Set^{\cK},\\
M^{\HH\PP}&\colon\HH\PP\cC\to\Set^{\cK},&
M^{\PP\HH}&\colon\PP\HH\cC\to\Set^{\cK},
\end{align*}
the first, third and fourth Heyting (Lemma~\ref{lem:PHeyt}), with
$M^{\HH\PP}\eta_{\PP\cC}\cong M^{\PP}$ and
\begin{equation}\label{eq:compat}
M^{\PP\HH}\circ c_\cC\;\cong\;M^{\HH\PP},
\end{equation}
both sides being Heyting functors whose restrictions to $\PP\cC$ are
coherent functors into a pretopos restricting to $M$ on $\cC$. The
functor $M^{\PP}$ is the Kripke model of $\PP\cC$ corresponding to $M$
under $\Mod(\PP\cC)\simeq\Mod(\cC)$: for $Y$ in $\PP\cC$, the set
$M^{\PP}(Y)(k)$ is the value at $Y$ of the extension of $M_k$, and an
arrow $w$ of $\cK$ acts by the extension of the homomorphism $M(w)$.

\begin{lem}[invariant]\label{lem:rigid}
Let $X$ be an object of $\HH\PP\cC$, $M$ a Kripke model of $\cC$ over
$\cK$, and $w\colon k\to l$ an arrow of $\cK$ such that $M(w)$ is an
isomorphism of models. Then $M^{\HH\PP}(X)(w)$ is injective.
\end{lem}
\begin{proof}
$X$ is a subobject of $\eta_{\PP\cC}Y$ for an object $Y$ of $\PP\cC$,
so $M^{\HH\PP}(X)$ is a subfunctor of $M^{\PP}(Y)$. The functor
$\cK\to\Mod(\PP\cC)$ corresponding to $M^{\PP}$ sends $w$ to an
isomorphism, since $M(w)$ is one and $\Mod(\PP\cC)\simeq\Mod(\cC)$; so
$M^{\PP}(Y)(w)$ is a bijection and its restriction is injective.
\end{proof}

When $\cC$ is not Boolean, some objects of $\PP\HH\cC$ do not have this
property (Section~5): they are quotients by relations whose truth at
$k$ depends on all the arrows out of $k$ and not only on the model
$M_k$.

\section{Non-Boolean categories have crossing homomorphisms}

\begin{lem}\label{lem:cross}
Let $S\le A$ be a subobject in a small coherent category $\cC$ which is
not complemented. Then there are models $M,N$ of $\cC$, a homomorphism
$h\colon M\to N$ and $a\in A(M)\setminus S(M)$ with $h_A(a)\in S(N)$.
\end{lem}
\begin{proof}
For $U\le B$ with inclusion $i$ we write $U(x)$ for the formula $\exists
u{:}U\,(i(u)=x)$. We use the completeness theorem: if $U,U_1,\dots,U_n\le B$ and
$U(M)\subseteq U_1(M)\cup\dots\cup U_n(M)$ for all models $M$, then
$U\le U_1\vee\dots\vee U_n$ \cite[D1.5]{Elephant}, \cite{MR77}; and the
compactness theorem for the classical first-order theory $T_\cC$ whose
models are the models of $\cC$, in many-sorted logic with possibly
empty sorts (the initial object is a sort).

Let $I=\{U\le A\mid U\wedge S=0\}$, an ideal of the distributive
lattice $\Sub(A)$, and $J$ the ideal of the $W$ with $W\le S\vee U$ for
some $U\in I$. It is proper: $A\le S\vee U$ with $U\wedge S=0$ would
make $U$ a complement of $S$. By the prime filter theorem there is a
prime filter $\mathfrak p$ of $\Sub(A)$ disjoint from $J$.

There are a model $M$ and $a\in A(M)$ with $a\in U(M)$ iff
$U\in\mathfrak p$, for all $U\le A$: the theory $T_\cC$ with a new
constant $c$ of sort $A$ and the axioms $U(c)$, $U\in\mathfrak p$, and
$\neg U(c)$, $U\notin\mathfrak p$, is finitely satisfiable, since
otherwise there are $U_0\in\mathfrak p$ and $U_1,\dots,U_n\notin
\mathfrak p$ with $U_0(M)\subseteq\bigcup_iU_i(M)$ in all models, hence
$U_0\le\bigvee_iU_i$, and $\bigvee_iU_i\in\mathfrak p$ contradicts the
primeness of $\mathfrak p$ (for $n=0$, $0\in\mathfrak p$ contradicts
its properness). In particular $a\notin S(M)$, and $a\notin U(M)$ for
$U\in I$.

Now add to $T_\cC$ a constant $c_b$ of sort $B$ for every object $B$
and every $b\in B(M)$, the axioms
\[
f(c_b)=c_{f_M(b)}\qquad(f\colon B\to B'\text{ a morphism of }\cC,\ b\in
B(M)),
\]
and the axiom $S(c_a)$. A finite part of these axioms mentions, besides
$c_a$, finitely many distinct constants $c_{b_1},\dots,c_{b_n}$,
$b_i\in B_i(M)$, and finitely many equations between them and $c_a$.
Put $b_0=a$, $B_0=A$ and $B=\prod_{i=0}^nB_i$, and let $\theta\le B$ be
the subobject defined by the conjunction of these equations, one
variable $y_i$ replacing each constant $c_{b_i}$ (an equation
$f(y_i)=y_j$ defines an equaliser, a subobject of $B$); then
$(a,b_1,\dots,b_n)\in\theta(M)$. If the finite part had no model, every
model $N'$ would satisfy $\exists_{\pi_0}\theta(N')\cap S(N')=
\emptyset$, so $U:=\exists_{\pi_0}\theta\le A$ would satisfy $U\wedge
S=0$ by completeness, that is $U\in I\subseteq J$; but $a\in U(M)$
gives $U\in\mathfrak p$, a contradiction. By compactness there is a
model $N$ of all these axioms, and $h_B(b):=c_b^N$ is, by the axioms, a
natural transformation $M\to N$, with $h_A(a)\in S(N)$.
\end{proof}

\begin{rem}
Lemma~\ref{lem:cross} is the homomorphism preservation theorem in the
form: a subobject whose complement is preserved by all homomorphisms of
models is complemented. Its converse is trivial, so a coherent category
is Boolean if and only if every homomorphism of its models reflects
every subobject. In that case every first-order formula is equivalent
to a coherent one and the homomorphisms are the elementary embeddings
of the models of the corresponding classical theory. (A single
homomorphism reflecting all subobjects need not be elementary: the
inclusion of a two-element set in a three-element set, for pure
equality.)
\end{rem}

\section{Proof of the main theorem}

\begin{proof}[Proof of Theorem~\ref{thm:main}]
The first assertion and (i)$\Leftrightarrow$(ii) are
Lemma~\ref{lem:ff}.

(iii)$\Rightarrow$(i). If $\cC$ is Boolean, so is $\PP\cC$
\cite[Prop.~3.4]{Morley}: complements are taken componentwise in
finite coproducts, and the complement of a saturated subobject is
saturated. By Lemma~\ref{lem:boolean}, $\eta_\cC$ and $\eta_{\PP\cC}$
are equivalences; so $\HH\PP\cC\simeq\PP\cC$ is a pretopos.

(i)$\Rightarrow$(iii). Suppose that $S\le A$ is not complemented in
$\cC$, and let $h\colon M\to N$ and $a$ be as in Lemma~\ref{lem:cross}.
In $\PP\HH\cC$ let $\neg S\le A$ be the pseudo-complement of (the image
of) $S$ in (the image of) $A$, and let
\[
Q\;:=\;(A+A)/E,\qquad E:=\Delta_{A+A}\vee(\kappa_1,\kappa_2)(\neg S)
\vee(\kappa_2,\kappa_1)(\neg S),
\]
with $\kappa_1,\kappa_2\colon A\to A+A$ the coprojections and
$(\kappa_i,\kappa_j)(\neg S)$ the image of $\neg S$ in $(A+A)\times
(A+A)$: the two copies of $A$ glued along $\neg S$, the pushout of
$A\leftarrow\neg S\to A$. The relation $E$ is an equivalence relation,
so $Q$ exists in the pretopos $\PP\HH\cC$. We show that $Q$ is not
isomorphic to $c_\cC(X)$ for any object $X$ of $\HH\PP\cC$.

Let $\cK$ be the fork, the poset with three elements $0<1$ and $0<2$,
with arrows $u\colon0\to1$ and $v\colon0\to2$, and $M^\ast$ the Kripke
model
\[
\begin{tikzcd}[column sep=large]
M^\ast_1=M&&M^\ast_2=N\\
&M^\ast_0=M\arrow[ul,"M^\ast(u)=\mathrm{id}"]\arrow[ur,"M^\ast(v)=h"']&
\end{tikzcd}
\]
Since $M^{\ast\PP\HH}$ is a Heyting functor between pretoposes, it
preserves pseudo-complements, coproducts, images and quotients; so
$M^{\ast\PP\HH}(Q)$ is the quotient of $M^\ast(A)+M^\ast(A)$ which
identifies the two copies of $x\in A(M^\ast_k)$ exactly when $x\in
(\neg M^\ast(S))(k)$. By \eqref{eq:neg}, $a\in(\neg M^\ast(S))(1)$,
because the only arrow out of $1$ is the identity and $a\notin S(M)$;
and $a\notin(\neg M^\ast(S))(0)$, because $M^\ast(A)(v)a=h_A(a)\in
S(N)$. Hence the two copies of $a$ are distinct points of
$M^{\ast\PP\HH}(Q)(0)$ with the same image under $u$ in
$M^{\ast\PP\HH}(Q)(1)$.

If $Q\cong c_\cC(X)$, then by \eqref{eq:compat} $M^{\ast\PP\HH}(Q)\cong
M^{\ast\HH\PP}(X)$, whose map along $u$ is injective by
Lemma~\ref{lem:rigid}, since $M^\ast(u)$ is an isomorphism. This is a
contradiction; so $c_\cC$ is not essentially surjective.

The last assertion is Lemma~\ref{lem:boolean} for $\cC$ and $\PP\cC$.
\end{proof}

\begin{cor}\label{cor:pretopos}
Let $\cC$ be a small pretopos. The free Heyting category $\HH\cC$ on
$\cC$ is a pretopos if and only if $\cC$ is Boolean, and then
$\HH\cC\simeq\cC$.
\end{cor}
\begin{proof}
$\iota\colon\cC\to\PP\cC$ is an equivalence, so $\HH\cC\simeq\HH\PP\cC$;
apply Theorem~\ref{thm:main}.
\end{proof}

\begin{ex}[pure equality]\label{ex:eq}
Let $\cC$ be the syntactic category of the coherent theory with one
sort $\sigma$ and no symbols and axioms. The diagonal of $\sigma\times
\sigma$ is not complemented; a crossing homomorphism is the map from a
two-element set $\{a,b\}$ to a one-element set, with the pair $(a,b)$.
The object $Q$ is two copies of $\sigma\times\sigma$ glued along
$\neg(x=y)$: an intuitionistic imaginary of the theory of pure
equality which is not an intuitionistically definable subset of a
coherent imaginary. Over the fork, $\neg(x=y)$ is empty at $0$ and is
the set of the two off-diagonal pairs at $1$; so the fibre of $Q$ at
$0$ has eight points, two copies of the four pairs, and the fibre at
$1$ has six, the two copies of $(a,b)$ and of $(b,a)$ being identified.
\end{ex}

\begin{ex}[propositional theories]\label{ex:prop}
Let $\cC=D$ be a bounded distributive lattice, a propositional coherent
theory. Up to isomorphism its models are the prime filters and the
homomorphisms are the inclusions, and Lemma~\ref{lem:cross} is the classical fact that a
distributive lattice in which the prime filters form an antichain is
Boolean. $\HH D$ is the free Heyting algebra on the distributive
lattice $D$. For the three-element chain $0<p<1$ the object $Q$ is
$(1+1)/E$ with the two points identified when $\neg p$; over the fork
$\{1\}=\{1\}\subset\{p,1\}$ of prime filters it has two points at $0$
and one at $1$, and no intuitionistically definable subobject of an
object of $\PP D$ behaves so.
\end{ex}

\section{Where the failure is: coproducts and quotients together}

The pretopos completion is $\PP=\EE\FF$, with $\FF$ the positive
completion (finite disjoint coproducts) and $\EE$ the exact completion
of a regular category \cite[A3.3]{Elephant}, \cite{TheoryComp}. Here a coherent category is \emph{positive} if it has finite
coproducts, disjoint and stable under pullback, $\FF$ is the reflector
onto the positive coherent categories (coherent functors from $\FF\cC$
to a positive coherent category correspond to those from $\cC$; the
unit is full and faithful and full on subobjects, and every object of
$\FF\cC$ is a finite coproduct of objects of $\cC$), and $\EE$ is the
exact (ex/reg) completion of a regular category. The object $Q$ uses a
coproduct and a quotient. The coproduct step alone is never
responsible; the quotient step alone is not responsible for
propositional theories, and is responsible under the hypothesis of
Proposition~\ref{prop:E}.

\begin{lem}\label{lem:FHeyt}
If $\cD$ is a Heyting category, then $\FF\cD$ is a Heyting category and
the unit $\cD\to\FF\cD$ is Heyting; and for a positive Heyting category
$\cD'$ a coherent functor $\FF\cD\to\cD'$ is Heyting if and only if its
restriction to $\cD$ is.
\end{lem}
\begin{proof}
Every object of $\FF\cD$ is a finite coproduct $\coprod_iA_i$ of
objects of $\cD$, the unit is full and faithful and full on subobjects,
and $\Sub(\coprod_iA_i)\cong\prod_i\Sub_\cD(A_i)$, $U\mapsto(U\wedge
A_i)_i$. Let $f\colon\coprod_iA_i\to\coprod_jB_j$, and let
$A_{ij}=A_i\wedge f^*B_j$, a subobject of $A_i$ in $\cD$, and
$f_{ij}\colon A_{ij}\to B_j$ the restriction, a morphism of $\cD$; for
each $i$ the $A_{ij}$ are disjoint with union $A_i$. For
$V=(V_j)_j$ one has $f^*V\wedge A_i=\bigvee_jf_{ij}^*V_j$, and
therefore, for $U\le\coprod_iA_i$,
\[
(\forall_fU)\wedge B_j\;=\;\bigwedge_i\forall_{f_{ij}}(U\wedge A_{ij}):
\]
$V\le\forall_fU$ iff $f_{ij}^*V_j\le U\wedge A_{ij}$ for all $i,j$ iff
$\bigvee_jf_{ij}^*V_j\le U\wedge A_i$ for all $i$, the $A_{ij}$ being
disjoint. The formula involves only the $\forall$ of $\cD$, finite
meets and the decomposition of coproducts; so the unit preserves
$\forall$, and so does a coherent functor to a positive Heyting
category whose restriction to $\cD$ is Heyting. The converse holds
since the unit is Heyting.
\end{proof}

\begin{prop}\label{prop:F}
For every small coherent category $\cC$ the canonical Heyting functor
$\HH\FF\cC\to\FF\HH\cC$ is an equivalence.
\end{prop}
\begin{proof}
By Lemma~\ref{lem:FHeyt}, $\FF\HH\cC$ is a positive Heyting category
with $\Heyt(\FF\HH\cC,\cD')\simeq\Heyt(\HH\cC,\cD')\simeq\Coh(\cC,
\cD')$ for positive Heyting categories $\cD'$. On the other side
$\Heyt(\HH\FF\cC,\cD')\simeq\Coh(\FF\cC,\cD')\simeq\Coh(\cC,\cD')$. So
it suffices to show that $\HH\FF\cC$ is positive; then both categories
are the free positive Heyting category on $\cC$, and the canonical
functor, which is compatible with the functors from $\cC$, is an
equivalence. Let $X\le\eta Y$ and $X'\le\eta Y'$ be objects of
$\HH\FF\cC$, with $Y,Y'$ in $\FF\cC$. The coherent functor $\eta$
sends the coproduct $Y+Y'$, the union of the disjoint subobjects $Y$
and $Y'$, to a union of disjoint subobjects, which in a coherent
category is a coproduct, disjoint and stable \cite[A1.4]{Elephant}; so
$X\vee X'\le\eta(Y+Y')$ is a disjoint and stable coproduct of $X$ and
$X'$, and the initial object is the empty one.
\end{proof}

\begin{cor}\label{cor:E}
Let $\cC$ be a small positive coherent category. Then $\HH\EE\cC$ is
positive, and it is exact if and only if $\cC$ is Boolean. For a
small bounded distributive lattice $D$, on the other hand, $\EE D\simeq D$ and
$\HH\EE D\simeq\HH D$ is exact, although $\HH\PP D$ is a pretopos only
if $D$ is Boolean.
\end{cor}
\begin{proof}
For positive $\cC$ the category $\EE\cC$ is a pretopos, equivalent to
$\PP\cC$ and to $\FF\EE\cC$; so $\HH\EE\cC\simeq\HH\FF\EE\cC$ is
positive by the proof of Proposition~\ref{prop:F}, and a positive
coherent category is a pretopos iff it is exact;
Theorem~\ref{thm:main} applies. In a preorder every equivalence
relation on an object $U$ is the whole of $U\times U\cong U$, the
kernel pair of the identity; so $D$ and the Heyting algebra $\HH D$ are
exact.
\end{proof}

For general coherent categories the quotient step alone does obstruct,
as soon as a model at which a subobject is crossed has two points. For
a coherent category $\cC$ the exact completion $\EE\cC$ is coherent and
effective (every equivalence relation is a kernel pair; a \emph{cover}
is a regular epimorphism), and restriction along the unit is an equivalence between the
coherent functors from $\EE\cC$ and from $\cC$ to an effective coherent
category, with all natural transformations \cite{TheoryComp}; in
particular $\Mod(\EE\cC)\simeq\Mod(\cC)$, and Lemma~\ref{lem:rigid}
holds for $\HH\EE\cC$ with the same proof: for an object $X$ of
$\HH\EE\cC$ and a Kripke model $M$ of $\cC$, the extension of $M$ to
$\HH\EE\cC$ sends $X$ to a functor whose maps along the arrows $w$ with
$M(w)$ an isomorphism are injective.

\begin{prop}\label{prop:E}
Let $\cC$ be a small coherent category, $S\le A$ a subobject,
$h\colon M\to N$ a homomorphism of models and $a\in A(M)\setminus S(M)$
with $h_A(a)\in S(N)$. If some object $B$ of $\cC$ has at least two
points in $M$, then $\HH\EE\cC$ is not effective: the equivalence
relation
\[
E\big((x,y),(x',y')\big)\;:=\;(x=x'\wedge y=y')\ \vee\ \big(x=x'\wedge
\neg S(x)\big)
\]
on $A\times B$ in $\HH\EE\cC$, which collapses the fibre over $x$ when
$\neg S(x)$, is not a kernel pair.
\end{prop}
\begin{proof}
$E$ is reflexive and symmetric, and transitive because each disjunct
contains $x=x'$ and the second depends only on $x$. Suppose that $E$ is
the kernel pair of a morphism of $\HH\EE\cC$; replacing it by the cover
in its image factorisation, $E$ is the kernel pair of a cover $q\colon
A\times B\to X$. Let $M^\ast$ be the Kripke model over the fork of the
proof of Theorem~\ref{thm:main} and $G$ its extension to $\HH\EE\cC$, a
Heyting functor to $\Set^{\cK}$. It preserves kernel pairs, covers and
pseudo-complements, so $G(X)(k)$ is the quotient of $A(M^\ast_k)\times
B(M^\ast_k)$ by the relation identifying $(x,y)$ and $(x,y')$ when
$x\in(\neg M^\ast(S))(k)$. As in that proof, $a\notin(\neg M^\ast(S))
(0)$ and $a\in(\neg M^\ast(S))(1)$; so for $y\ne y'$ in $B(M)$ the
points $(a,y)$ and $(a,y')$ are distinct in $G(X)(0)$ and have the same
image in $G(X)(1)$. This contradicts the invariant, $M^\ast(u)$ being
an isomorphism.
\end{proof}

\begin{ex}
For the theory of pure equality (Example~\ref{ex:eq}) take $A=\sigma
\times\sigma$, $S$ the diagonal, the map from a two-element set to a
one-element set, and $B=\sigma$: $\HH\EE\cC$ is not effective, so the Heyting completion does not
commute with the exact completion, without any coproduct. In a bounded
distributive lattice every object has at most one point in every
model, and Proposition~\ref{prop:E} does not apply, in accordance with
Corollary~\ref{cor:E}.
\end{ex}

So the Heyting completion commutes with adding coproducts; with adding
quotients it commutes for propositional theories and does not commute
when a subobject is crossed at a model in which some object has two
points. Whether it commutes with adding quotients in the remaining
cases, those of the non-Boolean categories in which every model at
which a subobject is crossed gives at most one point to every object,
is open (Question~\ref{que:E}).

\section{Rigid formulas are coherent, and the effective equivalence relations}

The invariant of Lemma~\ref{lem:rigid}, injectivity of the transition
maps of an object, corresponds for a subobject to the reflection of
membership along the same transitions (the two are related through
quotients, Theorem~\ref{thm:effective}). The coherent subobjects have
this property, and conversely; this
is a preservation theorem for intuitionistic logic, and it
determines which equivalence relations on coherent objects have
quotients in $\HH\PP\cC$.

\begin{defi}\label{def:rigid}
Let $U\le\eta A$ be a subobject in $\HH\cC$, $A$ an object of $\cC$.
For a Kripke model $M$ over $\cK$ write $U_M:=M^{\HH}(U)\subseteq M(A)$.
We say that $U$ is \emph{rigid} if for every Kripke model $M$, every
arrow $w\colon k\to l$ of $\cK$ such that $M(w)$ is an isomorphism of
models, and every $x\in A(M_k)$,
\[
M(w)_A(x)\in U_M(l)\ \Longrightarrow\ x\in U_M(k).
\]
\end{defi}

The converse implication always holds, $U_M$ being a subfunctor. For
every subobject $S\le A$ of $\cC$ the subobject $\eta S$ is rigid,
since $(\eta S)_M(k)=S(M_k)$.

\begin{lem}\label{lem:cosieve}
Let $\cK_0$ be a full subcategory of a small category $\cK$ such that
every arrow of $\cK$ with domain in $\cK_0$ has its codomain in
$\cK_0$. Then the restriction functor $\Set^{\cK}\to\Set^{\cK_0}$ is a
Heyting functor. Consequently, for a Kripke model $M$ over $\cK$ with
restriction $M_0$ to $\cK_0$, and $U\le\eta A$ in $\HH\cC$, one has
$U_{M_0}(k)=U_M(k)$ for every object $k$ of $\cK_0$.
\end{lem}
\begin{proof}
Restriction preserves limits and colimits, hence is coherent. For
$f\colon A\to B$ in $\Set^{\cK}$ and a subfunctor $U\subseteq A$,
\begin{multline*}
(\forall_fU)(k)=\{y\in B(k)\mid x\in U(l)\text{ for all }w\colon k\to
l\\\text{ and all }x\in A(l)\text{ with }f_l(x)=B(w)y\},
\end{multline*}
which involves only the arrows out of $k$; so restriction preserves
$\forall_f$. The restriction of $M^{\HH}$ is then a Heyting functor
extending $M_0$, hence isomorphic to $M_0^{\HH}$.
\end{proof}

We use the completeness theorem for Kripke models in the following
form, due to Joyal: for every small Heyting category $\cD$ there are a
small category $\cK$ and a conservative Heyting functor
$\cD\to\Set^{\cK}$ \cite{MR95}. For $\cD=\HH\cC$ such a functor is
$M^{\HH}$ for the Kripke model $M$ obtained by restriction to $\cC$; so
two subobjects $U,U'\le\eta A$ with $U_M=U'_M$ for all Kripke models
$M$ of $\cC$ are equal.

\begin{thm}[rigid formulas are coherent]\label{thm:rigid}
Let $\cC$ be a small coherent category and $U\le\eta A$ a subobject in
$\HH\cC$. Then $U$ is rigid if and only if $U=\eta S$ for a subobject
$S\le A$ of $\cC$.
\end{thm}
\begin{proof}
Let $U$ be rigid.

\emph{Step 1: $U_M(k)$ depends only on the model $M_k$.} Let $M$, $M'$
be Kripke models over $\cK$, $\cK'$, let $k$, $k'$ be objects of
$\cK$, $\cK'$ and $\theta\colon M_k\to M'_{k'}$ an isomorphism of
models; replacing $\cK'$ by a copy, we may assume that $\cK$ and
$\cK'$ are disjoint. Let $\cK''$ be the category whose objects are those of $\cK$,
those of $\cK'$ and a new object $o$, and whose arrows are those of
$\cK$ and $\cK'$, the identity of $o$, an arrow $\bar f\colon o\to j$
for every arrow $f\colon k\to j$ of $\cK$, and an arrow $\bar f'\colon
o\to j'$ for every arrow $f'\colon k'\to j'$ of $\cK'$, with
$g\bar f=\overline{gf}$. Let $M''$ be the Kripke model over $\cK''$
which is $M$ on $\cK$, $M'$ on $\cK'$, with $M''_o=M_k$,
$M''(\bar f)=M(f)$ and $M''(\bar f')=M'(f')\theta$. By
Lemma~\ref{lem:cosieve}, applied to $\cK$ and to $\cK'$, $U_{M''}$
agrees with $U_M$ on $\cK$ and with $U_{M'}$ on $\cK'$. Let $x\in
U_M(k)$. The arrow $\overline{\mathrm{id}_k}\colon o\to k$ acts by the
identity of $M_k$, so $x\in U_{M''}(o)$ by rigidity; and the arrow
$\overline{\mathrm{id}_{k'}}\colon o\to k'$ acts by $\theta$, so
$\theta_A(x)\in U_{M''}(k')=U_{M'}(k')$. By symmetry, $x\in U_M(k)$ iff
$\theta_A(x)\in U_{M'}(k')$.

\emph{Step 2: a class of models preserved by homomorphisms.} For a
model $N$ of $\cC$, a Kripke model over the terminal category, put
$U[N]:=U_N\subseteq A(N)$; it is the interpretation in $N$ of a
first-order formula $\varphi_U(x)$ of the language of the classical
theory $T_\cC$ of Lemma~\ref{lem:cross}, read classically. By Step 1,
$U_M(k)=U[M_k]$ for every Kripke model $M$ and every $k$. A
homomorphism $h\colon N\to N'$ is a Kripke model over the category
with one non-identity arrow, and $U_M$ is a subfunctor; so $h_A(U[N])
\subseteq U[N']$.

\emph{Step 3: preservation.} Let $\Phi$ be the set of the subobjects
$S\le A$ of $\cC$ with $S(N)\subseteq U[N]$ for all models $N$; it is
closed under finite unions. We claim that $U[N]\subseteq\bigcup_{S\in
\Phi}S(N)$ for all $N$. Let $a\in U[N]$, and consider the theory
$T_\cC$ with the constants $c_b$ and the axioms $f(c_b)=c_{f_N(b)}$ of
the proof of Lemma~\ref{lem:cross}, for the model $N$, and the axiom
$\neg\varphi_U(c_a)$. It has no model, since a model would give a
homomorphism $h\colon N\to N'$ with $h_A(a)\notin U[N']$. By
compactness a finite part has no model; as in that proof it gives a
subobject $\theta$ of a product $A\times\prod_iB_i$ and a point of
$\theta(N)$ over $a$ such that $S:=\exists_{\pi_0}\theta\le A$
satisfies $S(N')\subseteq U[N']$ for all models $N'$. So $S\in\Phi$
and $a\in S(N)$. By the claim and compactness, applied to the theory
$T_\cC$ with a constant $c$ and the axioms $\varphi_U(c)$ and $\neg
S(c)$, $S\in\Phi$, there is $S\in\Phi$ with $U[N]\subseteq S(N)$ for
all $N$, hence $U[N]=S(N)$ for all $N$.

\emph{Step 4.} For every Kripke model $M$ and every $k$,
$U_M(k)=U[M_k]=S(M_k)=(\eta S)_M(k)$. By the completeness theorem,
$U=\eta S$.
\end{proof}

Theorem~\ref{thm:rigid} is a preservation theorem: an intuitionistic
formula is equivalent, intuitionistically over the theory, to a
coherent formula if and only if its forcing at a node of a Kripke model
is reflected along the transitions which are isomorphisms.
Lemma~\ref{lem:cross} is the case $U=\neg\eta S$: the
pseudo-complement is rigid if and only if it is coherent, that is, if
and only if $S$ is complemented (if $\neg\eta S=\eta S'$, then in every
model $S'$ is the complement of $S$, the intuitionistic negation being
the complement in $\Set$, and $S'$ is a complement of $S$ in $\cC$ by
completeness).

Applied to the coherent category $\PP\cC$, whose Kripke models are
those of $\cC$, the theorem determines the effective equivalence
relations on the coherent objects of $\HH\PP\cC$.

\begin{thm}\label{thm:effective}
Let $\cC$ be a small coherent category, $Y$ an object of $\PP\cC$ and
$E$ an equivalence relation on $\eta Y$ in $\HH\PP\cC$. The following
are equivalent.
\begin{enumerate}
\item[(i)] $E$ is effective, that is, a kernel pair in $\HH\PP\cC$.
\item[(ii)] $E$ is rigid, as a subobject of $\eta(Y\times Y)$.
\item[(iii)] $E=\eta F$ for an equivalence relation $F$ on $Y$ in
$\PP\cC$.
\end{enumerate}
\end{thm}
\begin{proof}
(iii)$\Rightarrow$(i): $\eta$ is coherent, so $\eta F$ is the kernel
pair of $\eta(Y\to Y/F)$. (i)$\Rightarrow$(ii): let $E$ be the kernel
pair of a morphism, hence of a cover $q\colon\eta Y\to X$. For a
Kripke model $M$ the functor $M^{\HH\PP}$ preserves covers and kernel
pairs, so $M^{\HH\PP}(X)(k)$ is the quotient of $Y(M_k)$ by $E_M(k)$.
If $w\colon k\to l$ acts by an isomorphism and the images of $x,x'\in
Y(M_k)$ are related by $E_M(l)$, then the classes of $x$ and $x'$ have
the same image in $M^{\HH\PP}(X)(l)$, so they are equal by
Lemma~\ref{lem:rigid}: $(x,x')\in E_M(k)$. (ii)$\Rightarrow$(iii): by
Theorem~\ref{thm:rigid} for $\PP\cC$, $E=\eta F$ with $F\le Y\times Y$
in $\PP\cC$; and $F$ is an equivalence relation, because $\eta$
reflects the order of subobjects, the models of $\PP\cC$ being jointly
conservative and each of them being $N^{\HH}\eta$.
\end{proof}

So on the coherent imaginaries the free Heyting category has exactly
the quotients which were there before: no equivalence relation which
is new in intuitionistic logic has a quotient. This is analogous to
the modal quasi-pretoposes of \cite{GM25}, where the equivalence
relations required to be effective are those in a distinguished class
of subobjects; we do not verify any axioms of quasi-pretoposes for
$\HH\PP\cC$. For equivalence relations on an object $X\le\eta Y$ which
is not itself coherent, (i)$\Rightarrow$(ii) holds in the form that $E$
is reflected along the isomorphic transitions among the pairs of
points of $X$, with the same proof; the converse fails
(Proposition~\ref{prop:notsuff}).

\subsection*{The invariant does not characterise the image}
The invariant of Lemma~\ref{lem:rigid} compares a node with a node
above it which carries the same model. The objects in the image of
$c_\cC$ satisfy a stronger condition, in which the models at several
nodes are compared through homomorphisms which need not be transitions
of the Kripke model.

\begin{lem}[cocone invariant]\label{lem:factor}
Let $X\le\eta_{\PP\cC}Y$ be an object of $\HH\PP\cC$, $Y$ in $\PP\cC$,
let $M$ be a Kripke model of $\cC$ over $\cK$, $J\colon\mathcal J\to
\cK$ a functor from a small category (for instance the inclusion of a
subcategory $\cK'$), $N$ a model of $\cC$, and $g_j\colon M_{J(j)}\to
N$, for the objects $j$ of $\mathcal J$, homomorphisms with $g_{j'}
\circ M(J(w))=g_j$ for every arrow $w\colon j\to j'$ of $\mathcal J$.
Write $\bar g_j\colon Y(M_{J(j)})\to Y(N)$ for the extension of $g_j$
to $\PP\cC$ at $Y$. Let $j_0,\dots,j_n$ be objects of $\mathcal J$ and
$x_i\in M^{\HH\PP}(X)(J(j_i))$ points such that for every $i<n$ there
is an arrow of $\mathcal J$ between $j_i$ and $j_{i+1}$, in one of the
two directions, whose image under $J$ sends one of $x_i$, $x_{i+1}$ to
the other; we say that $x_0$ and $x_n$ are \emph{connected by a zigzag
in $\mathcal J$}. Then $\bar g_{j_0}(x_0)=\bar g_{j_n}(x_n)$. In
particular, if $J(j_0)=J(j_n)$ and $g_{j_0}=g_{j_n}$ is an
isomorphism, then $x_0=x_n$.
\end{lem}
\begin{proof}
$M^{\HH\PP}(X)$ is a subfunctor of $M^{\PP}(Y)$, whose value at $k$ is
$Y(M_k)$ and whose map along an arrow is the extension of the
transition. The extension of homomorphisms to $\PP\cC$ is functorial;
so $\bar g_{j'}\circ M^{\PP}(Y)(J(w))=\bar g_j$ for the arrows
$w\colon j\to j'$ of $\mathcal J$, and $\bar g_{j_i}(x_i)$ does not
depend on $i$. The extension of an isomorphism is a bijection.
\end{proof}

Lemma~\ref{lem:rigid} is the case of the category $\mathcal J$ with one
non-identity arrow $j\to j'$, sent to $w\colon k\to l$ with $M(w)$ an
isomorphism, $N=M_k$, $g_j=\mathrm{id}$, $g_{j'}=M(w)^{-1}$ and the
points $x,\ M^{\HH\PP}(X)(w)x=M^{\HH\PP}(X)(w)x',\ x'$. Another case: if
$u\colon r\to k$ and $v\colon r\to l$ are arrows with $M(u)=g\circ
M(v)$ for a homomorphism $g\colon M_l\to M_k$, then two points at $r$
with the same image under $v$ have the same image under $u$ (take the
arrow $v$, $N=M_k$, $g_r=M(u)$, $g_l=g$). For a subcategory $\cK'$ we
speak of a zigzag in $\cK'$.

\begin{prop}\label{prop:notsuff}
Let $D$ be the three-element chain $0<p<1$, so that $\HH D$ is the free
Heyting algebra on the generator $p$. Let
\[
a:=\neg\neg p\Rightarrow p,\quad A:=a\Rightarrow(p\vee\neg p),\quad
B:=A\Rightarrow(\neg\neg p\vee a),\quad W:=\neg\neg p,
\]
so that $W\le A\wedge B$, and let $Z:=(A+B)/E$ in $\PP\HH D$ be the
pushout of $A\leftarrow W\to B$, the subobjects $A$ and $B$ of $1$
glued along $W$: with $\kappa_A$, $\kappa_B$ the coprojections,
$E:=\Delta_{A+B}\vee(\kappa_A,\kappa_B)(W)\vee(\kappa_B,\kappa_A)(W)$.
\begin{enumerate}
\item[(a)] For every Kripke model $M$ of $D$ over a small category,
all the transition maps of $M^{\PP\HH}(Z)$ are injective. In particular
$Z$ has the invariant of Lemma~\ref{lem:rigid}.
\item[(b)] $Z$ is not isomorphic to $c_D(X)$ for any object $X$ of
$\HH\PP D$.
\item[(c)] The same formulas define an equivalence relation on an
object $A+B\le\eta(1+1)$ of $\HH\PP D$ which is reflected along all
transitions and is not effective.
\end{enumerate}
\end{prop}
\begin{proof}
The models of $D$ are the prime filters $P_0=\{1\}\subset P_1=\{p,1\}$.
So a Kripke model over $\cK$ is a functor $k\mapsto|k|$ from $\cK$ to
the chain $0<1$; a subobject of $1$ in $\Set^{\cK}$ is a set of objects
closed under arrows, and we write $k\Vdash U$ when $k$ belongs to the
value of $U\in\HH D$; thus $k\Vdash p$ iff $|k|=1$. From $\neg\neg p$
and $a$ one gets $p$; so $W\le A$, and $W\le B$ is clear. The fibre of
$Z$ at $k$ consists of a point $s$ if $k\Vdash A$ and a point $t$ if
$k\Vdash B$, and $s=t$ there iff $k\Vdash\neg\neg p$.

(a) $A\wedge B\le\neg\neg p\vee a$ and $a\wedge A\le p\vee\neg p\le
\neg\neg p\vee\neg p$; so $A\wedge B\le\neg\neg p\vee\neg p$. If the
fibre at $k$ has two points, then $k\Vdash A\wedge B$ and
$k\nVdash\neg\neg p$, so $k\Vdash\neg p$; then no $l$ with an arrow
$k\to l$ forces $\neg\neg p$, and the two points have distinct images
at $l$.

(b) Let $M$ be the Kripke model over the poset with elements
$0,\dots,5$ and order generated by the arrows shown, with $|0|=1$ and
$|k|=0$ for $k\ne0$:
\[
\begin{tikzcd}[column sep=large,row sep=normal]
0&1\\
2\arrow[u]&3\arrow[u]\arrow[ul]\\
4\arrow[u]\arrow[uur,bend left=10]&5\arrow[u]\arrow[ul]
\end{tikzcd}
\]
(the arrow from $4$ goes to $1$). The nodes above $1$, $2$, $3$, $4$,
$5$ are $\{1\}$, $\{2,0\}$, $\{3,1,0\}$, $\{4,2,1,0\}$,
$\{5,3,2,1,0\}$. One finds: $0\Vdash p$; $1\Vdash\neg p$, hence
$1\Vdash a$, $A$, $B$ and $1\nVdash\neg\neg p$; $2\Vdash\neg\neg p$,
hence $2\Vdash A$, $B$, and $2\nVdash a$; $3\Vdash a$, since
$\neg\neg p$ is forced only at $0$ among the nodes above $3$, and
$3\nVdash p\vee\neg p$, so $3\nVdash A$; $4\Vdash A$, since among the
nodes above $4$ only $1$ and $0$ force $a$, and they force $p\vee\neg
p$; $4\nVdash\neg\neg p\vee a$, so $4\nVdash B$; $5\nVdash A$, since
$3$ is above $5$; and $5\Vdash B$, since the nodes above $5$ which
force $A$ are $2$, $1$, $0$, and they force $\neg\neg p\vee a$. So the
fibres of $M^{\PP\HH}(Z)$ at $4$, $2$, $5$, $1$ are $\{s\}$, $\{s=t\}$,
$\{t\}$, $\{s,t\}$ with $s\ne t$ at $1$.

Suppose $Z\cong c_D(X)$; then $M^{\PP\HH}(Z)\cong M^{\HH\PP}(X)$ by
\eqref{eq:compat}. Let $\cK'$ be the full subposet on $1,2,4,5$; all
its nodes carry the model $P_0$, with identities as transitions, and
$g_k=\mathrm{id}_{P_0}$ is a cocone. The points
\[
s\text{ at }1,\quad s\text{ at }4,\quad s=t\text{ at }2,\quad
t\text{ at }5,\quad t\text{ at }1
\]
are related successively by the arrows $4\to1$, $4\to2$, $5\to2$,
$5\to1$. By Lemma~\ref{lem:factor}, $s=t$ at $1$; a contradiction.

(c) We use the same letters for the subobjects of $\eta1$ in
$\HH\PP D$ defined by the same formulas. In $\HH\PP D$ let $A+B$ be the
subobject of $\eta(1+1)$ which is $A$ on the first summand and $B$ on
the second; the formula for $E$ defines an equivalence relation $E_0$
on it with $c_D(E_0)=E$. It is reflected along all transitions by (a).
If $E_0$ were the kernel pair of a morphism, it would be the kernel
pair of a cover $A+B\to X$, and $c_D$, a Heyting functor, preserves
covers and kernel pairs; so $c_D(X)\cong Z$, contradicting (b).
\end{proof}

So injectivity of the transition maps, along the isomorphisms of
models or even along all transitions, does not characterise the
essential image of $c_\cC$, already for the smallest non-Boolean
propositional theory: $Z$ is the union of two subobjects $A$ and $B$ which lie in the
image, all its transition maps are injective, and it is excluded only
by a zigzag through four nodes carrying the same model. For two copies
of a coherent object $\eta Y$ glued along $W$, on the other hand, the
invariant makes $W$ rigid, hence coherent (Theorem~\ref{thm:rigid}),
and the glued object is coherent.

\subsection*{The image for finite propositional theories}
For a finite propositional theory the cocone invariant characterises
the image, and it has to be tested in one finite Kripke model only.

\begin{lem}\label{lem:subimage}
For every small coherent category $\cC$, an object of $\PP\HH\cC$ is in
the essential image of $c_\cC$ if and only if it admits a monomorphism
to $(\PP\eta)(Y)$ for some object $Y$ of $\PP\cC$.
\end{lem}
\begin{proof}
An object of $\HH\PP\cC$ is a subobject $X\le\eta_{\PP\cC}Y$, and
$c_\cC(\eta_{\PP\cC}Y)=(\PP\eta)(Y)$. Conversely let $V\le(\PP\eta)(Y)$,
and let $q\colon\coprod_jA_j\to Y$ be a cover in $\PP\cC$ with the
$A_j$ in $\cC$. The pullback of $V$ along $(\PP\eta)(q)$ is a family of
subobjects $V_j\le\eta A_j$, which belong to $\HH\cC$, the unit
$\HH\cC\to\PP\HH\cC$ being full on subobjects \cite[A3.3]{Elephant}.
Let $\HH\iota\colon\HH\cC\to\HH\PP\cC$ be the Heyting functor induced
by the unit $\iota\colon\cC\to\PP\cC$; then $c_\cC\circ\HH\iota$ is
isomorphic to the unit $\HH\cC\to\PP\HH\cC$, both being Heyting
functors which restrict to $(\PP\eta)\iota$ on $\cC$. In $\HH\PP\cC$
let $V'\le\eta_{\PP\cC}Y$ be the image under $\eta_{\PP\cC}(q)$ of the
union of the subobjects $(\HH\iota)(V_j)$ of the summands of
$\eta_{\PP\cC}(\coprod_jA_j)$. Then $c_\cC(V')$ is the image of
$\coprod_jV_j$ under the cover $(\PP\eta)(q)$, which is $V$.
\end{proof}

Let now $D$ be a finite distributive lattice and $X$ the poset of its
prime filters, the models of $D$; the homomorphisms are the
inclusions, and $\PP D$ is the category $\mathbf{FinSet}^X$ of functors
from $X$ to finite sets, the coherent objects of the classifying topos
$\Set^X$ \cite[D3.3]{Elephant}; in the same way $\PP D'=
\mathbf{FinSet}^{X'}$ for every finite distributive lattice $D'$ with
poset of prime filters $X'$, an element of $D'$ corresponding to the
subterminal functor supported on the prime filters containing it. Let
$H_0\subseteq H_1\subseteq\cdots$ be the sublattices of $\HH D$ with
$H_0=D$ and $H_{n+1}$ generated by $D$ and the elements $u\Rightarrow
v$ with $u,v\in H_n$. They are finite, being finitely generated
distributive lattices, and their union is $\HH D$. Let $X_n$ be the
poset of the prime filters of $H_n$, ordered by inclusion, and
$|F|:=F\cap D$; so $F\mapsto|F|$ is a Kripke model of $D$ over $X_n$.

\begin{lem}\label{lem:truth}
In the Kripke model $X_n$, for $U\in H_n$ and $F\in X_n$: $F\Vdash U$
if and only if $U\in F$.
\end{lem}
\begin{proof}
By induction on $k\le n$ we prove it for $U\in H_k$. The filters being
prime and forcing commuting with finite meets and joins, it suffices to
treat the generators. For $d\in D$, $F\Vdash d$ iff $d\in|F|$. Let
$u,v\in H_{k-1}$. If $u\Rightarrow v\in F$ and $F\subseteq F'\Vdash u$,
then $u\in F'$ by induction, so $v\in F'$ and $F'\Vdash v$; hence
$F\Vdash u\Rightarrow v$. If $u\Rightarrow v\notin F$, let $\hat F$ be a
prime filter of $\HH D$ with $\hat F\cap H_n=F$; it exists because the
filter of $\HH D$ generated by $F$ and the ideal generated by
$H_n\setminus F$ are disjoint. The filter generated by $\hat F$ and $u$
does not contain $v$, since $f\wedge u\le v$ with $f\in\hat F$ gives
$u\Rightarrow v\in\hat F\cap H_n$. So there is a prime filter $\hat G
\supseteq\hat F$ of $\HH D$ with $u\in\hat G$, $v\notin\hat G$, and
$F':=\hat G\cap H_n\supseteq F$ has $F'\Vdash u$ and $F'\nVdash v$ by
induction. Hence $F\nVdash u\Rightarrow v$.
\end{proof}

Every object of $\PP\HH D$ is a quotient $Z=(\coprod_{i\in I}A_i)/E$ of
a finite coproduct of elements $A_i$ of $\HH D$ by an equivalence
relation, which is given by elements $E_{ij}\le A_i\wedge A_j$; the
fibre of $Z$ at a node $k$ of a Kripke model is the set of the $i$ with
$k\Vdash A_i$ modulo the relation $k\Vdash E_{ij}$.

\begin{thm}\label{thm:image}
Let $D$ be a finite distributive lattice, $Z=(\coprod_iA_i)/E$ an
object of $\PP\HH D$, and $n$ such that all $A_i$ and $E_{ij}$ belong
to $H_n$. The following are equivalent.
\begin{enumerate}
\item[(i)] $Z$ is in the essential image of $c_D$.
\item[(ii)] For every Kripke model $M$ of $D$ over a poset $\cK$ and
every node $k$: two points of the fibre of $Z$ at $k$ which are
connected by a zigzag, in the sense of Lemma~\ref{lem:factor}, through
nodes $l$ with $M_l\subseteq M_k$, are equal.
\item[(iii)] Condition (ii) for the finite Kripke model $X_n$.
\end{enumerate}
If they hold, $Z$ is a subobject of $(\PP\eta)(Y)$ for the functor $Y$
on $X$ with $Y(x)=\mathrm{colim}_{|F|\subseteq x}Z(F)$, the left Kan
extension along $F\mapsto|F|$ of the fibres of $Z$ on $X_n$.
\end{thm}
\begin{proof}
(i)$\Rightarrow$(ii): if $Z\cong c_D(X')$, then $M^{\PP\HH}(Z)\cong
M^{\HH\PP}(X')$, and Lemma~\ref{lem:factor} applies to the full
subposet $\cK'$ of the nodes $l$ with $M_l\subseteq M_k$, the model
$N=M_k$ and the inclusions $g_l\colon M_l\subseteq M_k$, with
$g_k=\mathrm{id}$. (ii)$\Rightarrow$(iii) is clear.

(iii)$\Rightarrow$(i). The inclusions $D\subseteq H_n\subseteq\HH D$
are homomorphisms of lattices and induce coherent functors
\[
\PP D=\mathbf{FinSet}^X\longrightarrow\PP H_n=\mathbf{FinSet}^{X_n}
\longrightarrow\PP\HH D
\]
with composite $\PP\eta$, the first being composition with
$F\mapsto|F|$. The elements $E_{ij}$ define an equivalence relation on
$\coprod_iA_i$ in $\PP H_n$, the conditions being inequalities between
elements of $H_n$ which hold in $\HH D$; let $Z_n$ be the quotient, an
object of $\PP H_n$ whose image in $\PP\HH D$ is $Z$. As a functor on
$X_n$, $Z_n(F)$ is the set of the $i$ with $A_i\in F$ modulo the
relation $E_{ij}\in F$; by Lemma~\ref{lem:truth} this is the fibre of
$Z$ at $F$ in the Kripke model $X_n$. Let $Y$ be the left Kan extension
of $Z_n$ along $F\mapsto|F|$, a functor from $X$ to finite sets: $Y(x)$
is the colimit of $Z_n(F)$ over the $F$ with $|F|\subseteq x$, the set
of the pairs $(F,z)$, $z\in Z_n(F)$, modulo the equivalence relation
generated by $(F,z)\sim(F',Z_n(F\subseteq F')z)$. The unit
$Z_n(F)\to Y(|F|)$, $z\mapsto[(F,z)]$, is natural in $F$, and it is
injective if and only if two points of $Z_n(F)$ connected by a zigzag
through filters $F'$ with $|F'|\subseteq|F|$ are equal, which is
(iii). So the unit is a monomorphism from $Z_n$ to the image of $Y$ in
$\PP H_n$; applying the coherent functor to $\PP\HH D$ gives a
monomorphism $Z\to(\PP\eta)(Y)$, and Lemma~\ref{lem:subimage} applies.
\end{proof}

So for a finite propositional theory the essential image of $c_D$
consists of the objects with the cocone invariant, and the embedding
into a coherent object, when it exists, is the canonical one computed
in the finite model $X_n$. Given the finite lattice $D$ and formulas for the $A_i$ and
$E_{ij}$, of implication depth at most $n$, the lattice $H_n$ and the
poset $X_n$ can be computed, intuitionistic propositional logic being
decidable; so membership in the image is decidable. The object $Z$ of
Proposition~\ref{prop:notsuff} fails (ii) in the six-element model
given there. Theorem~\ref{thm:image} uses the finiteness of $D$ through
the finite models $X_n$, in which every element of $H_n$ is forced
exactly where it should be; the same method applies to the first-order
theories with finitely many finite models considered in
Theorem~\ref{thm:presheaf}.

\subsection*{Infinite propositional theories}
For an infinite propositional theory the image is still described
through finite Kripke models, those of the finite subtheories; but the
cocone invariant in the Kripke models of the theory itself is no longer
sufficient. Let $D$ be a bounded distributive lattice. A finite
sublattice $D''$ (containing $0$ and $1$) gives a homomorphism
$\HH D''\to\HH D$ and coherent functors $\PP D''\to\PP D$ and
$\PP\HH D''\to\PP\HH D$, compatible with $\PP\eta$. A morphism of
$\PP\HH D$ from $(\coprod_iA_i)/E$ to $(\coprod_jB_j)/F$ is a
functional relation, a family of elements $R_{ij}\le A_i\wedge B_j$ of
$\HH D$ with
\[
A_i\le\textstyle\bigvee_jR_{ij},\qquad R_{ij}\wedge R_{ij'}\le
F_{jj'},\qquad E_{i'i}\wedge R_{ij}\wedge F_{jj'}\le R_{i'j'},
\]
and it is a monomorphism iff $R_{ij}\wedge R_{i'j}\le E_{ii'}$.

\begin{cor}\label{cor:allD}
Let $D$ be a bounded distributive lattice and $Z=(\coprod_iA_i)/E$ an
object of $\PP\HH D$, the $A_i$ and $E_{ij}$ being given by formulas in
the elements of a finite sublattice $D'$. Then $Z$ is in the essential
image of $c_D$ if and only if there is a finite sublattice $D''\supseteq
D'$ of $D$ such that the same formulas define an object $Z_{D''}$ of
$\PP\HH D''$ which is in the essential image of $c_{D''}$, that is,
which satisfies the conditions of Theorem~\ref{thm:image} for $D''$.
\end{cor}
\begin{proof}
If $Z_{D''}$ is a subobject of $(\PP\eta)(Y)$ with $Y$ in $\PP D''$,
its image $Z$ in $\PP\HH D$ is a subobject of $(\PP\eta)(Y')$, $Y'$ the
image of $Y$ in $\PP D$; Lemma~\ref{lem:subimage} applies. Conversely,
let $Z\to(\PP\eta)(Y)$ be a monomorphism, $Y=(\coprod_jd_j)/F$ with
$d_j$, $F_{jl}$ in $D$, given by elements $R_{ij}$ of $\HH D$. The
free Heyting algebra functor is a left adjoint, and directed colimits
of Heyting algebras and of lattices are computed on the underlying
sets; so $\HH D$ is the directed colimit of the $\HH D''$, $D''$ the
finite sublattices of $D$. The elements $R_{ij}$ are given by formulas
in finitely many elements of $D$, and the finitely many inequalities
which say that $E$ is an equivalence relation and that $R$ is a
monomorphism $Z\to(\PP\eta)(Y)$ hold in $\HH D$, hence in $\HH D''$ for
some finite $D''$ containing $D'$, the $d_j$, the $F_{jl}$ and these
elements. So $Z_{D''}$ is defined and is a subobject of
$(\PP\eta)(Y_{D''})$.
\end{proof}

\begin{prop}\label{prop:infinite}
There are a countable bounded distributive lattice $D$ and an object
$Z$ of $\PP\HH D$ which is not in the essential image of $c_D$,
although it has the cocone invariant: for every Kripke model $M$ of $D$
over a small category and every node $k$, two points of the fibre of
$Z$ at $k$ which are connected by a zigzag through nodes $l$ with
$M_l\subseteq M_k$ are equal.
\end{prop}
\begin{proof}
Let $X$ be the set of the points $\top$, $\infty$, $y_i$, $w_i$
($i\ge0$) and $g_i$, $h_i$ ($i\ge1$), ordered by $g_i<y_{i-1}$,
$g_i<y_i$, $h_i<w_{i-1}$, $h_i<w_i$ and $x<\top$ for all $x\ne\top$:
two fences, and the point $\infty$, below a top element.
\[
\begin{tikzcd}[column sep=tiny]
y_0&&y_1&&y_2&\cdots&\infty&\cdots&w_2&&w_1&&w_0\\
&g_1\arrow[ul]\arrow[ur]&&g_2\arrow[ul]\arrow[ur]&&&&&&h_2\arrow[ul]
\arrow[ur]&&h_1\arrow[ul]\arrow[ur]&
\end{tikzcd}
\]
Give $X$ the topology in which every point other than $\infty$ is
isolated and the neighbourhoods of $\infty$ are the cofinite sets
containing it; $X$ is compact and Hausdorff. Call the index of
$y_i,w_i,g_i,h_i$ the number $i$. If $x\not\le x'$ there is a clopen
up-set containing $x$ and not $x'$: for $x\ne\infty$ the finite set
$\uparrow\!x$, and for $x=\infty$ the set $U_N$ of $\top$, $\infty$,
the $y_j,w_j$ with $j\ge N$ and the $g_j,h_j$ with $j>N$, for $N$
larger than the index of $x'$. So $X$ is a Priestley space; let $D$ be
the lattice of its clopen up-sets, which is countable. By Priestley
duality the prime filters of $D$ are the filters $\{U\mid x\in U\}$,
$x\in X$, and inclusion of prime filters is the order of $X$. So a
Kripke model of $D$ is a functor $\nu\mapsto|\nu|$ to the poset $X$.
Let
\[
t:=\{\top\},\ \ a:=\{y_0,\top\},\ \ b:=\{w_0,\top\},\quad
A:=a\vee\neg t,\ \ W:=b\wedge\neg t\le A,
\]
and let $Z$ be the two copies of $A$ glued along $W$. The fibre of $Z$
at a node $\nu$ has two points $s_\nu,t_\nu$ if $\nu\Vdash A$, they are
equal iff $\nu\Vdash W$, and the transitions preserve the two names.

\emph{The cocone invariant.} Suppose that $s_k\ne t_k$ are connected by
a zigzag through nodes $l$ with $|l|\le|k|$. As the transitions
preserve the names, the zigzag contains a node forcing $W$; let
$k=\nu_0,\dots,\nu_j=m$ be its part up to the first such node. All
$\nu_i$ force $A$, and consecutive ones are joined by an arrow. Let
$F$ be the set of the points $w_i$, $h_i$, and say that a node $\nu$
is \emph{good} if $\nu\Vdash\neg t$ and $|\nu|\in F$. The node $m$ is
good: it forces $b$ and $\neg t$, so $|m|\in\{w_0,\top\}$ and
$|m|\ne\top$. If $\nu_{i+1}$ is good, so is $\nu_i$: if there is an
arrow $\nu_{i+1}\to\nu_i$, then $\nu_i\Vdash\neg t$, so $|\nu_i|\ne
\top$, and $|\nu_i|\ge|\nu_{i+1}|$, and the points above a point of $F$
are in $F\cup\{\top\}$; if there is an arrow $\nu_i\to\nu_{i+1}$, then
$|\nu_i|\le|\nu_{i+1}|$ is in $F$, so $\nu_i\nVdash a$, and
$\nu_i\Vdash A$ gives $\nu_i\Vdash\neg t$. Hence $k$ is good, and
$w_0=|m|\le|k|\in F$ gives $|k|=w_0$. So $k\Vdash b\wedge\neg t=W$,
contradicting $s_k\ne t_k$.

\emph{$Z$ is not in the image.} By Corollary~\ref{cor:allD} it
suffices to show that for every finite sublattice $D''$ containing
$t,a,b$ the object $Z_{D''}$ violates condition (ii) of
Theorem~\ref{thm:image}. A clopen set either contains $\infty$ and is
cofinite, or does not contain $\infty$ and is finite. So there is
$N\ge1$ such that every element of $D''$ either contains $\infty$ and
all points of index $\ge N$, or contains none of them. Write
$x'':=\{U\in D''\mid x\in U\}$, a prime filter of $D''$; then
$x\mapsto x''$ is monotone, and $x''=\infty''$ for all $x$ of index
$\ge N$. Let $\cK$ be the finite poset with elements $k$, $r$, $e$,
$\bar y_i,\bar w_i$ ($0\le i<N$), $\bar g_i,\bar h_i$ ($1\le i\le N$)
and order generated by
\[
r<k,\quad r<\bar y_0,\quad\bar g_i<\bar y_{i-1},\ \bar g_i<\bar y_i,
\quad\bar h_i<\bar w_{i-1},\ \bar h_i<\bar w_i,
\]
where $\bar y_N:=\bar w_N:=e$, and $M$ the Kripke model of $D''$ with
$|k|=\top''$, $|r|=y_0''$, $|e|=\infty''$, $|\bar y_i|=y_i''$,
$|\bar g_i|=g_i''$, $|\bar w_i|=w_i''$, $|\bar h_i|=h_i''$; it is
monotone because $g_N''=h_N''=\infty''$. The elements $t,a,b$ belong to
$x''$ exactly for $x=\top$, for $x\in\{y_0,\top\}$ and for
$x\in\{w_0,\top\}$. So $k$ is the only node forcing $t$, every node
other than $k$ and $r$ forces $\neg t$, and $k$ and $r$ force $a$; all
nodes force $A$, and $\bar w_0$ is the only node forcing $W$. Hence
$s_k\ne t_k$; and $s_k$ and $t_k$ are connected by the zigzag
\[
k\leftarrow r\to\bar y_0\leftarrow\bar g_1\to\bar y_1\leftarrow\cdots
\to e\leftarrow\bar h_N\to\cdots\to\bar w_1\leftarrow\bar h_1\to
\bar w_0
\]
with the points $s$, followed by the same nodes in the reverse order
with the points $t$, since $s=t$ at $\bar w_0$. All its nodes $l$
satisfy $|l|\subseteq|k|=\top''$, every non-empty up-set containing
$\top$.
\end{proof}

In the space $X$ of Proposition~\ref{prop:infinite} the two fences
are not connected in $X\setminus\{\top\}$, but their images are
connected, without the image of $\top$ and through the image of
$\infty$, in the poset of prime filters of every finite sublattice
containing $t,a,b$, by zigzags for whose length there is no bound independent of the
sublattice. So for infinite theories the invariant has to be read in the
models of the finite subtheories, or to be supplemented by a condition
of continuity; this is the form in which the description of the image
has to be sought for first-order theories (Question~\ref{que:image}).

\subsection*{Finite theories of presheaf type}
Theorem~\ref{thm:image} is the propositional case of a statement about
first-order theories with finitely many finite models. Let $\mathcal X$
be a finite category in which idempotents split, and $\mathcal P:=
\mathbf{FinSet}^{\mathcal X}$, a pretopos, the category of coherent
objects of the presheaf topos $\Set^{\mathcal X}$
\cite[D3.3]{Elephant}. Its models are, up to isomorphism, the evaluations $A\mapsto A(m)$
at the objects $m$ of $\mathcal X$, and the homomorphisms between
these are the arrows of $\mathcal X$ (the points of a presheaf topos are the flat
functors \cite[B3.2]{Elephant}, and a flat functor on a finite category
with split idempotents is representable). So a Kripke model of
$\mathcal P$ over $\cK$ is a functor $\cK\to\mathcal X$. For a finite
poset $\mathcal X=X$ this is the situation of Theorem~\ref{thm:image},
with $\PP D=\mathcal P$; for a finite groupoid $\mathcal P$ is Boolean;
for the category $\mathcal X$ obtained from a finite monoid by
splitting its idempotents, $\mathcal P$ is the category of the finite
sets with an action of the monoid. Let $\cC$ be a small coherent category
with $\PP\cC\simeq\mathcal P$. By the proof of Lemma~\ref{lem:ff},
$\PP\HH\cC\simeq\PP\HH\mathcal P$, and $c_\cC$ corresponds to the unit
$\HH\mathcal P\to\PP\HH\mathcal P$; so we may assume $\cC=\mathcal P$,
and the question is which objects of $\PP\HH\mathcal P$ are isomorphic
to objects of $\HH\mathcal P$.

Write $y(x)=\mathcal X(x,-)$, and $y(\gamma)\colon y(x')\to y(x)$ for
an arrow $\gamma\colon x\to x'$. Every $A$ in $\mathcal P$ is covered
by the morphisms $q_a\colon y(x)\to A$, $a\in A(x)$, with
$q_a(\varphi)=A(\varphi)(a)$. We omit $\eta$ in the notation for
subobjects in $\HH\mathcal P$. Define finite sublattices
$\mathcal L_n(x)\subseteq\Sub_{\HH\mathcal P}(y(x))$: $\mathcal L_0(x)=
\Sub_{\mathcal P}(y(x))$, and $\mathcal L_{n+1}(x)$ is generated by
$\mathcal L_n(x)$ and the elements $u\Rightarrow v$ ($u,v\in
\mathcal L_n(x)$), $y(\gamma)^*u$ ($\gamma\colon x'\to x$, $u\in
\mathcal L_n(x')$), $\exists_{y(\gamma)}u$ and $\forall_{y(\gamma)}u$
($\gamma\colon x\to x'$, $u\in\mathcal L_n(x')$). For $A$ in
$\mathcal P$ let $\mathcal L_n(A)$ be the set of the $U\le A$ in
$\HH\mathcal P$ with $q_a^*U\in\mathcal L_n(x)$ for all $x$ and
$a\in A(x)$.

\begin{lem}\label{lem:levels}
\begin{enumerate}
\item[(a)] $\mathcal L_n(A)$ is a sublattice, $f^*\mathcal L_n(B)
\subseteq\mathcal L_n(A)$ for $f\colon A\to B$ in $\mathcal P$, and
$u\Rightarrow v$, $\exists_fu$, $\forall_fu$ are in level $n+1$ for
$u,v$ in level $n$. In particular $\mathcal L_n(y(x))=\mathcal L_n(x)$.
\item[(b)] $\Sub_{\HH\mathcal P}(A)=\bigcup_n\mathcal L_n(A)$ for every
$A$ in $\mathcal P$.
\end{enumerate}
\end{lem}
\begin{proof}
(a) First, $y(\gamma)^*\mathcal L_n(x')\subseteq\mathcal L_n(x)$ for
$\gamma\colon x'\to x$, by induction on $n$: pullback commutes with the
lattice operations and with $\Rightarrow$, and for $\delta\colon x'\to
x_1$ and $u\in\mathcal L_{n-1}(x_1)$ the pullback of
$\exists_{y(\delta)}u$ along $y(\gamma)\colon y(x)\to y(x')$ is
computed on the pullback $P$ of $y(\delta)$ and $y(\gamma)$ in
$\mathcal P$; covering $P$ by representables $p_i\colon y(x_i)\to P$,
whose composites with the projections to $y(x)$ and $y(x_1)$ are
$y(\gamma_i)$ and $y(\delta_i)$, we get
$y(\gamma)^*\exists_{y(\delta)}u=\bigvee_i\exists_{y(\gamma_i)}
y(\delta_i)^*u$ and $y(\gamma)^*\forall_{y(\delta)}u=\bigwedge_i
\forall_{y(\gamma_i)}y(\delta_i)^*u$, by the Beck--Chevalley conditions
of a Heyting category and because the $p_i$ are jointly covering; and
$y(\delta_i)^*u\in\mathcal L_{n-1}(x_i)$ by induction. This
gives $\mathcal L_n(y(x))=\mathcal L_n(x)$, since $q_\gamma=y(\gamma)$
for $\gamma\in y(x)(x')$. The other assertions follow in the same way:
$q_a^*f^*U=q_{f(a)}^*U$; and for $f\colon A\to B$, $b\in B(x)$ and the
pullback $P$ of $f$ and $q_b$, covered by representables $p_i$ with
composites $y(\gamma_i)\colon y(x_i)\to y(x)$ and $q_{a_i}\colon
y(x_i)\to A$, one has $q_b^*\exists_fU=\bigvee_i\exists_{y(\gamma_i)}
q_{a_i}^*U$ and $q_b^*\forall_fU=\bigwedge_i\forall_{y(\gamma_i)}
q_{a_i}^*U$.

(b) $\HH\mathcal P$ is the syntactic category of the internal theory of
$\mathcal P$ in intuitionistic logic, with the objects of $\mathcal P$
as sorts; so the subobjects of the objects $A$ of $\mathcal P$ in
$\HH\mathcal P$ are obtained from the subobjects in $\mathcal P$, which
are in level $0$, by the operations of (a).
\end{proof}

Let $\cK_n$ be the following finite category. An object is a pair
$(m,G)$, $m$ an object of $\mathcal X$ and $G$ a family of subsets
$G(U)\subseteq\mathcal X(x,m)$, $U\in\mathcal L_n(x)$, of the form
$G(U)=\hat G(U)$ for a coherent functor $\hat G\colon\HH\mathcal P\to
\Set$ whose restriction to $\mathcal P$ is the evaluation at $m$. An
arrow $(m,G)\to(m',G')$ is an arrow $\varphi\colon m\to m'$ with
$\varphi\circ G(U)\subseteq G'(U)$ for all $U$. For $U\in\mathcal L_n
(A)$ put $G(U):=\{a\in A(m)\mid\mathrm{id}_m\in G(q_a^*U)\}=\hat G(U)$;
the arrows of $\cK_n$ preserve these sets, by Lemma~\ref{lem:levels}(a).
The functor $(m,G)\mapsto m$ is a Kripke model of $\mathcal P$ over
$\cK_n$.

\begin{lem}\label{lem:truthP}
In the Kripke model $\cK_n$, for $A$ in $\mathcal P$ and $U\in
\mathcal L_n(A)$, the value of $U$ at the node $(m,G)$ is $G(U)$.
\end{lem}
\begin{proof}
Both sides are compatible with pullback along the $q_a$, with finite
meets and joins and with images; so by induction on $k\le n$ it
suffices to treat $u\Rightarrow v$ and $\forall_fu$, $f=y(\gamma)
\colon y(x')\to y(x)$, with $u,v$ in level $k-1$. Let $\alpha\in
\mathcal X(x,m)$. If $\alpha\in G(\forall_fu)$ and $\varphi\colon(m,G)
\to(m',G')$, $\beta\in\mathcal X(x',m')$ with $f(\beta)=\varphi\alpha$,
then $\varphi\alpha\in G'(\forall_fu)$ and $\beta\in\hat G'(f^*\forall_f
u)\subseteq\hat G'(u)$, which is the value of $u$ at $(m',G')$ by
induction; so $\alpha$ is in the value of $\forall_fu$. Conversely let
$\alpha\notin\hat G(\forall_fu)$. As in the proof of
Lemma~\ref{lem:cross}, consider the classical theory of the coherent
functors on $\HH\mathcal P$ with constants $c_b$ for the elements $b$
of the sets $\hat G(B)$, the axioms $g(c_b)=c_{g(b)}$ for the morphisms
$g$ of $\HH\mathcal P$, a new constant $d$ of sort $y(x')$ and the
axioms $f(d)=c_\alpha$ and $\neg u(d)$. If it had no model, a finite
part would give a subobject $V\le y(x)$ in $\HH\mathcal P$ with
$\alpha\in\hat G(V)$ and $f^*V\le u$, hence $V\le\forall_fu$,
contradicting $\alpha\notin\hat G(\forall_fu)$. A model is a coherent
functor $\hat G'$ with a natural transformation $\hat G\to\hat G'$ and
an element $\beta\in\hat G'(y(x'))\setminus\hat G'(u)$ over the image
of $\alpha$. The restriction of $\hat G'$ to $\mathcal P$ is the
evaluation at some $m'$, the natural transformation restricts to an
arrow $\varphi\colon m\to m'$, which is an arrow $(m,G)\to(m',G')$ of
$\cK_n$, and $\beta\notin G'(u)$; so $\alpha$ is not in the value of
$\forall_fu$. The case of $u\Rightarrow v$ is the same, with the
axioms $u(c_\alpha)$ and $\neg v(c_\alpha)$.
\end{proof}

\begin{lem}\label{lem:definable}
Let $A$ be in $\mathcal P$ and $S$ a subfunctor of $(m,G)\mapsto A(m)$
on $\cK_n$. Then there is $V\le A$ in $\HH\mathcal P$ with $\hat G(V)=
S(m,G)$ for every coherent functor $\hat G$ on $\HH\mathcal P$ inducing
the object $(m,G)$.
\end{lem}
\begin{proof}
For an object $(m,G)$ let $\tau_G\in\mathcal L_n(m)$ be the meet of the
$U\in\mathcal L_n(m)$ with $\mathrm{id}_m\in G(U)$. For an object
$(m',G')$ and $\varphi\colon m\to m'$ we have $\varphi\in G'(\tau_G)$
if and only if $\varphi$ is an arrow $(m,G)\to(m',G')$: if it is,
$\mathrm{id}_m\in G(\tau_G)$ gives $\varphi\in G'(\tau_G)$; conversely,
for $\alpha\in G(U)$, $U\in\mathcal L_n(x)$, we have $\mathrm{id}_m\in
G(y(\alpha)^*U)$, so $\tau_G\le y(\alpha)^*U$ and $\varphi\alpha\in
G'(U)$. Let $V$ be the join of the $\exists_{q_a}\tau_G$ over the
objects $(m,G)$ and the $a\in S(m,G)$. Then $\hat G'(\exists_{q_a}
\tau_G)$ is the set of the $A(\varphi)(a)$, $\varphi\colon(m,G)\to
(m',G')$, which is contained in $S(m',G')$, $S$ being a subfunctor, and
contains $a$ if $(m',G')=(m,G)$.
\end{proof}

\begin{thm}\label{thm:presheaf}
Let $\mathcal X$ be a finite category in which idempotents split,
$\mathcal P=\mathbf{FinSet}^{\mathcal X}$, and $Z=(\coprod_iU_i)/E$ an
object of $\PP\HH\mathcal P$, with $U_i\le A_i$ and $E_{ij}\le A_i
\times A_j$ in level $n$, the $A_i$ in $\mathcal P$. The following are
equivalent.
\begin{enumerate}
\item[(i)] $Z$ is isomorphic to an object of $\HH\mathcal P$.
\item[(ii)] For every Kripke model $M\colon\cK\to\mathcal X$ and every
node $k$: two points of the fibre of $Z$ at $k$ are equal if they are
connected by a zigzag in the category $\mathcal J$ of the pairs
$(l,g)$, $l$ a node and $g\colon M_l\to M_k$ an arrow of $\mathcal X$,
whose arrows $(l,g)\to(l',g')$ are the arrows $w\colon l\to l'$ with
$g'M(w)=g$, starting and ending at $(k,\mathrm{id})$.
\item[(iii)] Condition (ii) for the finite Kripke model $\cK_n$.
\end{enumerate}
\end{thm}
\begin{proof}
(i)$\Rightarrow$(ii) is Lemma~\ref{lem:factor} for the functor
$\mathcal J\to\cK$, $(l,g)\mapsto l$, the model $N=M_k$ and the cocone
$g_{(l,g)}=g$. (ii)$\Rightarrow$(iii) is clear. (iii)$\Rightarrow$(i):
by Lemma~\ref{lem:truthP} the fibre $Z(m,G)$ of $Z$ at a node of
$\cK_n$ is the set of the pairs $(i,a)$, $a\in G(U_i)$, modulo the
relation $(a,a')\in G(E_{ij})$. Let $Y$ be the left Kan extension of
$(m,G)\mapsto Z(m,G)$ along $(m,G)\mapsto m$: $Y(m)$ is the colimit of
$Z(l,G)$ over the finite category of the objects $(l,G)$ of $\cK_n$
with an arrow $g\colon l\to m$, a finite set; so $Y$ is an object of
$\mathcal P$. The unit $\theta\colon Z(m,G)\to Y(m)$ is natural on
$\cK_n$, and it is injective by (iii). For each $i$ the sets
$\{(a,\theta[i,a])\mid a\in G(U_i)\}\subseteq(A_i\times Y)(m)$ form a
subfunctor on $\cK_n$; let $V_i\le A_i\times Y$ be as in
Lemma~\ref{lem:definable}. That the $V_i$ form a functional relation
from $Z$ to $Y$ which is a monomorphism is expressed by finitely many
inequalities between subobjects, in $\HH\mathcal P$, of objects of
$\mathcal P$, built from the $U_i$, $E_{ij}$, $V_i$ by the operations
of a coherent category. Every coherent functor $\hat G$ on
$\HH\mathcal P$ induces an object $(m,G)$ of $\cK_n$, with $\hat G(V_i)$
the graph of $\theta$ on the $i$-th summand; so these inequalities
hold under every $\hat G$, and by the completeness theorem for the
coherent category $\HH\mathcal P$ they hold in $\HH\mathcal P$. So
there is a monomorphism $Z\to(\PP\eta)(Y)$ in $\PP\HH\mathcal P$, and
Lemma~\ref{lem:subimage} applies.
\end{proof}

For a finite poset $\mathcal X$ conditions (i) and (ii) are those of
Theorem~\ref{thm:image}: an arrow $M_l\to M_k$ exists iff
$M_l\subseteq M_k$, and it is unique. In general the zigzags pass through
homomorphisms between the models at the nodes which need not be
inclusions, and through automorphisms; for a finite group $\mathcal X$
every object of $\PP\HH\mathcal P=\mathcal P$ satisfies the conditions.

\begin{ex}[a twisted form]\label{ex:twist}
Let $\mathcal X$ be the product of the poset $0<1$ and the group
$\mathbb Z/2$, a category with two objects in which idempotents split.
An object of $\mathcal P$ is an equivariant map $A(0)\to A(1)$ of
finite $\mathbb Z/2$-sets. Let $p\le1$ be the subterminal object with
$p(0)=\emptyset$ and $p(1)$ a point, and $\sigma:=y(0)$, the free orbit at both objects;
the complement $x\ne y$ of the diagonal of $\sigma\times\sigma$ is a
subobject in $\mathcal P$. So $\mathcal P$ contains the three-element
chain $0<p<1$ and a sort with exactly two elements and decidable
equality, and a Kripke model over a poset $\cK$ is a functor
$k\mapsto|k|$ to $0<1$ together with a two-element set $\sigma_k$ at
every node and bijections along the arrows. Let $a$, $A$, $B$ be the
formulas of Proposition~\ref{prop:notsuff}; recall that $\neg\neg p\le
A\wedge B\le\neg\neg p\vee\neg p$. So $A\wedge B$ is the union of the
disjoint subobjects
\[
W_1:=\neg\neg p,\qquad W_2:=\neg p\wedge A\wedge B.
\]
Let $Z$ be the quotient of $A\times\sigma+B\times\sigma$ which
identifies $x$ in the first summand with $x$ in the second over $W_1$,
and with the other element $y\ne x$ of $\sigma$ over $W_2$. This is an
equivalence relation, because $\sigma$ has exactly two elements and
decidable equality. The fibre of $Z$ at a node $k$ is $\sigma_k$ if $k$
forces only one of $A$, $B$, and two copies of $\sigma_k$ identified,
straight or crosswise, if $k\Vdash A\wedge B$; all its transition maps
are injective, since the forcing of $\neg\neg p$ and of $\neg p$ is
preserved along the transitions and the two are disjoint. So $Z$ is a
twisted form of $(A\vee B)\times\sigma$, which is the object obtained
by identifying straight over $W_1$ and $W_2$, and which lies in
$\HH\mathcal P$. But $Z$ is not isomorphic to an object of
$\HH\mathcal P$. Take the six-element Kripke model of
Proposition~\ref{prop:notsuff}, with $\sigma_k=\{u,v\}$ at every node
and identities as transitions. With the indices $s$, $t$ for the two
summands, the fibres at $4$ and $5$ are $\{u_s,v_s\}$ and
$\{u_t,v_t\}$; at $2$, which forces $W_1$, $u_s=u_t$ and $v_s=v_t$; and
at $1$, which forces $W_2$, $u_s=v_t$ and $v_s=u_t$. The points
\[
u_s\text{ at }1,\quad u_s\text{ at }4,\quad u_s=u_t\text{ at }2,\quad
u_t\text{ at }5,\quad u_t=v_s\text{ at }1
\]
form a zigzag along $4\to1$, $4\to2$, $5\to2$, $5\to1$ through nodes
with the same model and identities as transitions, from $u_s$ to the
distinct point $v_s$ of the fibre at $1$. This contradicts condition
(ii), with the homomorphisms $g=\mathrm{id}$. The propositional part
of this example is harmless: the quotient of $Z$ by the action of
$\mathbb Z/2$ on $\sigma$ is the subterminal object $A\vee B$. The
obstruction comes from the automorphism of $\sigma$: $Z$ is isomorphic
to $A\times\sigma$ over $A$ and to $B\times\sigma$ over $B$, and the
two identifications differ by the exchange of the two elements over
$W_2$ and agree over $W_1$.
\end{ex}

\section{When excluded middle repairs the comparison}

Let $\KK\colon\Heyt\to\mathbf{BCoh}$ be the classicalisation, the
reflector onto the full sub-$2$-category $\mathbf{BCoh}$ of Boolean
coherent categories which adds excluded
middle to the internal first-order theory; $\BB\simeq\KK\HH$ is the
Morleyization, and $\KK\PP\cD\simeq\PP\KK\cD$ for every Heyting
category $\cD$ \cite[Thm.~9.2]{Morley}. The comparison
$b_\cC\colon\BB\PP\cC\to\PP\BB\cC$ of \cite{Morley} is $\KK c_\cC$
followed by this equivalence at $\cD=\HH\cC$
\cite[Cor.~9.3]{Morley}:
\[
\begin{tikzcd}[column sep=large]
\HH\PP\cC\arrow[r,"c_\cC"]\arrow[d]&\PP\HH\cC\arrow[d]\\
\BB\PP\cC=\KK\HH\PP\cC\arrow[r,"\KK c_\cC"]\arrow[rr,bend right=18,
"b_\cC"']&\KK\PP\HH\cC\arrow[r,"\simeq"]&\PP\KK\HH\cC=\PP\BB\cC.
\end{tikzcd}
\]

\begin{cor}\label{cor:K}
For a small coherent category $\cC$ the following are equivalent: the
Morleyization commutes with the pretopos completion at $\cC$, that is
$b_\cC$ is an equivalence; the Galois criterion of
\cite[Thm.~6.1]{Morley} holds for $\cC$ (in the saturated, strongly
homogeneous models of the completions of the classical theory, every
element of a classical imaginary sort has the same automorphism
stabiliser as some element of a coherent imaginary sort of the same
model); the Heyting functor $c_\cC$ is
inverted by $\KK$. If these conditions hold and $\cC$ is not Boolean, then $c_\cC$ is a
full and faithful Heyting functor which is not an equivalence and
becomes one when excluded middle is added. This is the case for the
theory of pure equality, for the theory of an equivalence relation,
for the theories without axioms over a monadic relational signature
with at least one sort, and for the non-Boolean propositional theories
\cite[Ex.~6.5, Thm.~8.2, Thm.~7.1, Prop.~5.2]{Morley}.
\end{cor}
\begin{proof}
The first equivalence is \cite[Thm.~6.1]{Morley}, the second the
factorisation above, and the rest is Theorem~\ref{thm:main}. The
theories listed are not Boolean: in each there is a homomorphism of
models which is not an elementary embedding (a non-injective map of
sets in the first and third cases, a map merging two classes, a proper
inclusion of prime filters).
\end{proof}

In the example of pure equality one sees the repair: classically
$\neg(x=y)$ is a complemented subobject of $\sigma\times\sigma$ and the
glued object $Q$ of Example~\ref{ex:eq} becomes the coproduct of
$\sigma\times\sigma$ and the diagonal, a classical subobject of a
coherent imaginary. For the theory of complete multipartite graphs
\cite[Thm.~5.3]{Morley} the repair fails: $\KK$ does not invert
$c_\cC$. Thus the statement of \cite{Morley} that the forced order of
the Morleyization and the completion arises in the Heyting completion
is to be read as follows. The comparison is never an equivalence
before excluded middle is added, except trivially; the Galois criterion
decides whether excluded middle repairs it.

\section{Questions}

\begin{que}\label{que:E}
Let $\cC$ be a non-Boolean coherent category such that, whenever a
homomorphism $h\colon M\to N$ crosses a subobject, every object of
$\cC$ has at most one point in $M$ (for instance the product of a
non-Boolean bounded distributive lattice and a Boolean coherent
category). Is
$\HH\EE\cC$ effective?
\end{que}

\begin{que}\label{que:image}
Describe the essential image of $c_\cC$ in $\PP\HH\cC$ for first-order
theories with infinite models, for instance for the theory of pure
equality. The invariant of Lemma~\ref{lem:rigid} is necessary and not
sufficient (Proposition~\ref{prop:notsuff}); the cocone invariant of
Lemma~\ref{lem:factor} is necessary, it is sufficient for the finite
theories of presheaf type (Theorems~\ref{thm:image}
and~\ref{thm:presheaf}) and not for infinite propositional theories
(Proposition~\ref{prop:infinite}), where the image is described through
the finite subtheories (Corollary~\ref{cor:allD}).
\end{que}

\begin{que}
Let $L$ be an intermediate predicate logic, given by axiom schemes
added to intuitionistic first-order logic, and $\HH_L$ the
corresponding translation. The fork is not a frame for the logic of linear frames
(G\"odel--Dummett), and over a finite linear order two consecutive
nodes whose transition is an isomorphism of models force the same
formulas, under this isomorphism; so the proof of
Theorem~\ref{thm:main} does not apply. For which $L$ does $\HH_L$
commute with $\PP$ beyond the Boolean case?
\end{que}

\begin{que}
Is there a description of the class of coherent functors inverted by
$\HH$, and of its relation to the classes $V_\BB$ and $V_\PP$ of
\cite{Morley} in the lattice of two-out-of-three classes, in spite of
$\HH$ not being a reflector?
\end{que}

\section*{Declarations}
The author has no competing interests to declare.

\begin{thebibliography}{99}
\bibitem{Elephant} P.~T.~Johnstone, \emph{Sketches of an Elephant: A Topos
Theory Compendium}, Oxford Univ. Press, 2002.
\bibitem{GM25} S.~Ghilardi, J.~Marqu\`es, \emph{First-order modal logic
via logical categories}, preprint, arXiv:2504.02985v1 (2025).
\bibitem{MR95} M.~Makkai, G.~E.~Reyes, \emph{Completeness results for
intuitionistic and modal logic in a categorical setting}, Ann. Pure
Appl. Logic 72 (1995), 25--101.
\bibitem{MR77} M.~Makkai, G.~E.~Reyes, \emph{First Order Categorical
Logic}, Lecture Notes in Math. 611, Springer, 1977.
\bibitem{Morley} R.~Shibusawa, \emph{Translations in the lattice of
two-out-of-three classes: Morleyization as a sameness of theories, its
join with Morita equivalence, and the forced order of translation and
completion}, preprint (2026), archived at DOI 10.5281/zenodo.24122457.
\bibitem{TheoryComp} R.~Shibusawa, \emph{Comparing theories in the lattice
of two-out-of-three classes: Morita equivalence as a join of completions,
the forced order of coproducts and quotients, and conceptual completeness
as an identity of samenesses}, preprint (2026), archived at DOI
10.5281/zenodo.24095158.
\end{thebibliography}

\end{document}
